\documentclass[11pt]{article}

\usepackage[a4paper,margin=2.5cm]{geometry}
\usepackage{booktabs}
\usepackage{amsmath}
\usepackage{hyperref}

\hypersetup{
  colorlinks=true,
  linkcolor=black,
  citecolor=black,
  urlcolor=blue,
  pdftitle={ODD Protocol Description of the Dark Patterns Trust-Erosion ABM}
}

\title{Supplementary Material: ODD Protocol Description of the\\
Dark Patterns Trust-Erosion Agent-Based Model}

\author{%
Dejana Pivac\thanks{Faculty of Informatics, Juraj Dobrila University of Pula, Pula, Croatia. \texttt{dpivac@fipu.hr}}
\and
Darko Etinger\thanks{Faculty of Informatics, Juraj Dobrila University of Pula, Pula, Croatia. \texttt{darko.etinger@unipu.hr}}
}

\date{2026}

\newcommand{\code}[1]{\texttt{#1}}

\begin{document}

\maketitle

\begin{abstract}
\noindent This document provides a complete description of the Dark Patterns
Trust-Erosion agent-based model (ABM) following the ``Overview, Design concepts,
Details'' (ODD) protocol, in its updated form as specified by
Grimm et al. (2020). The model simulates how the deployment of manipulative
interface design (``dark patterns'') by a single digital platform erodes user
trust, propagates negative word-of-mouth across a social network, and ultimately
drives churn and reputational decline over a multi-year horizon. The description
below is grounded directly in the model source code (the Python package
\code{backend/app/simulation/}); every quantitative parameter is reported with
the name of the constant that defines it in \code{config.py} so that the
description and the implementation remain in one-to-one correspondence.
All numerical \emph{results} produced by the model (means, confidence intervals,
trigger fractions) are reported in the main manuscript and its tables/figures and
are regenerated by the reproduction harness \code{reproduce.py}; this protocol
documents only the model's \emph{structure and mechanisms}.
\end{abstract}

\tableofcontents

\bigskip

This protocol follows the seven-element ODD structure of
Grimm et al. (2020): (1) Purpose and patterns; (2) Entities, state variables,
and scales; (3) Process overview and scheduling; (4) Design concepts;
(5) Initialization; (6) Input data; and (7) Submodels.

\section{Purpose and Patterns}
\label{sec:purpose}

\subsection{Purpose}

The purpose of the model is to study, as an emergent dynamic process, how a
digital platform's use of dark patterns degrades the trust relationship between
the platform and a heterogeneous population of users, and how that degradation
translates, through individual trust erosion, social transmission, and
exit decisions, into platform-level outcomes (churn, reputation loss, and a
gap between short-term extractive revenue and forgone long-term revenue). The
model is explanatory rather than predictive: it is designed to expose the
qualitative regimes and threshold behaviours that arise when manipulative design
is intensified, and to compare those regimes across dark-pattern types and user
populations. It backs a thesis and an accompanying journal manuscript on the
societal consequences of dark patterns.

\subsection{Patterns}

In the ODD sense, ``patterns'' are the characteristic, multi-scale phenomena the
model is built to reproduce and that serve as informal criteria for its
structural realism. This model is designed to reproduce four such patterns,
implemented as explicitly tracked \emph{tipping points} (see
Section~\ref{sec:tipping} for thresholds):

\begin{enumerate}
  \item \textbf{Trust Collapse}: the gradual, then accelerating, decline of
  mean user trust toward a low equilibrium when dark-pattern intensity is high.
  \item \textbf{Social Contagion}: a self-reinforcing rise in negative
  word-of-mouth (WOM) as harmed users warn their network neighbours.
  \item \textbf{Churn Cascade}: the conversion of trust loss and social
  contagion into a large cumulative fraction of users abandoning the platform.
  \item \textbf{Extractive Divergence}: the widening gap between short-term
  revenue (boosted by dark-pattern extraction) and discounted long-term revenue,
  co-occurring with sustained churn; the signature of a platform that is
  ``cashing out'' its user base.
\end{enumerate}

A further qualitative pattern the model is expected to reproduce is
\textbf{population heterogeneity in vulnerability}: distinct user types (skeptic,
naive, activist) should exhibit different trust trajectories, churn timing, and
WOM behaviour, with naive users retained longest (they rationalise bad
experiences) and activists/skeptics reacting fastest.

\section{Entities, State Variables, and Scales}
\label{sec:entities}

\subsection{Entities}

The model contains the following entity types:

\begin{itemize}
  \item \textbf{User agents} (\code{UserAgent}): heterogeneous consumers, one
  per network node. The default population is $N = 500$ (\code{DEFAULT\_NUM\_AGENTS}).
  Each user is assigned a discrete \emph{user type} (skeptic, naive, or activist)
  that fixes the ranges from which its individual traits are drawn.
  \item \textbf{Platform agent} (\code{PlatformAgent}): a single platform
  operator that deploys dark patterns at a chosen global intensity, provides
  customer support, and optionally adapts its strategy.
  \item \textbf{Dark-pattern objects} (\code{DarkPattern}): three immutable
  descriptors (\code{forced\_trial}, \code{hard\_cancel}, \code{drip\_pricing}),
  each carrying detectability, harm, and revenue parameters.
  \item \textbf{The model / environment} (\code{DarkPatternTrustModel}):
  owns the social network, the scheduling, and all model-level economic and
  reputational state.
  \item \textbf{The social network}: a graph whose nodes are user agents and
  whose edges carry word-of-mouth.
\end{itemize}

\subsection{User-agent state variables}

User traits are \emph{static} attributes sampled once at initialization from
type-specific ranges (\code{USER\_TYPE\_RANGES}); the dynamic variables evolve
during the run. Tables~\ref{tab:traits} and~\ref{tab:dynamic} list them.
Trait values are drawn by \code{beta\_sample(rng, low, high)}, a symmetric
$\mathrm{Beta}(5,5)$ draw rescaled to $[\text{low}, \text{high}]$
(\code{BETA\_SHAPE}~$=5.0$; bell-shaped, extremes rare).

\begin{table}[!ht]
\centering
\caption{Static user-agent traits and their sampling ranges by user type
(\code{USER\_TYPE\_RANGES} in \code{config.py}). Each value is a
$\mathrm{Beta}(5,5)$ draw scaled to the listed interval. The
\code{pattern\_sensitivity} range is applied independently to each of the three
dark patterns.}
\label{tab:traits}
\small
\begin{tabular}{lccc}
\toprule
Trait & Skeptic & Naive & Activist \\
\midrule
\code{trust\_baseline}            & $[0.55, 0.70]$ & $[0.75, 0.90]$ & $[0.60, 0.75]$ \\
\code{digital\_literacy}          & $[0.70, 0.95]$ & $[0.15, 0.40]$ & $[0.75, 0.95]$ \\
\code{manipulation\_sensitivity}  & $[0.55, 0.75]$ & $[0.30, 0.50]$ & $[0.70, 0.90]$ \\
\code{social\_activity}           & $[0.30, 0.50]$ & $[0.10, 0.30]$ & $[0.45, 0.70]$ \\
\code{complaint\_propensity}      & $[0.40, 0.65]$ & $[0.10, 0.25]$ & $[0.40, 0.60]$ \\
\code{switching\_cost}            & $[0.30, 0.45]$ & $[0.50, 0.70]$ & $[0.25, 0.40]$ \\
\code{trust\_resilience}          & $[0.00, 0.10]$ & $[0.30, 0.50]$ & $[0.00, 0.05]$ \\
\code{pattern\_sensitivity}       & $[1.0, 1.4]$   & $[0.6, 0.9]$   & $[1.1, 1.5]$   \\
\bottomrule
\end{tabular}
\end{table}

\begin{table}[!ht]
\centering
\caption{Dynamic user-agent state variables. All trust-like quantities are
clamped to $[0,1]$. ``Initial'' indicates the value at step 0.}
\label{tab:dynamic}
\small
\begin{tabular}{lll}
\toprule
Variable & Meaning & Initial \\
\midrule
\code{trust}                 & current trust in the platform                  & $=$\,\code{trust\_baseline} \\
\code{perceived\_fairness}   & perceived fairness of the platform             & $=$\,\code{trust\_baseline} \\
\code{harm}                  & accumulated cumulative harm (saturating)       & $0.0$ \\
\code{negative\_wom}         & current negative-WOM propensity                & $0.0$ \\
\code{warning\_awareness}    & primed suspicion (raises detection)            & $0.0$ \\
\code{positive\_sentiment}   & primed positive sentiment                      & $0.0$ \\
\code{cumulative\_exposure}  & sum of intensities of exposures received       & $0.0$ \\
\code{exposure\_count}       & number of exposures received                   & $0$ \\
\code{pattern\_exposure\_count} & per-pattern exposure counter (buildup ramp) & $0$ \\
\code{active}                & membership flag (\code{False} once churned)    & \code{True} \\
\code{received\_negative\_signal} & social signal staged for this step        & $0.0$ \\
WOM counters                 & \code{negative\_wom\_sent/received}, \code{positive\_wom\_sent/received} & $0$ \\
\bottomrule
\end{tabular}
\end{table}

\subsection{Platform state variables}

The \code{PlatformAgent} holds \code{dark\_pattern\_intensity}~$\in[0,1]$
(a global scalar applied to every active pattern),
\code{customer\_support\_quality}~$\in[0,1]$ (trust-recovery effectiveness),
the boolean \code{adaptive\_platform}, an internal \code{reputation}~$\in[0,1]$
(a trust/WOM aggregate), and \code{short\_term\_revenue} /
\code{long\_term\_revenue} accumulators.

\subsection{Model-level (environmental) state variables}

The model object additionally holds a separate economic and reputational layer:
\code{platform\_reputation} (a $0$--$100$ score, distinct from the agent's
$[0,1]$ \code{reputation}); per-step and cumulative revenue, costs, profit, and
net value; a no-dark-pattern \code{opportunity\_cost} projection;
\code{cumulative\_churn} and \code{churn\_rate}; per-pattern step counters for
exposures, detections, and trust loss; the four tipping-point status records
with their consecutive-step streak counters; and the random number generator.

\subsection{Scales}

\begin{itemize}
  \item \textbf{Temporal:} one model step represents \textbf{one week}. The
  reproduction horizon is \textbf{312 steps} ($\approx$ six years;
  \code{DEFAULT\_MAX\_STEPS}~$=312$).
  \item \textbf{Spatial / network:} there is no physical space. Interaction
  topology is a \textbf{Watts--Strogatz small-world} network
  (\code{DEFAULT\_NETWORK\_TYPE}~$=$ \code{"small\_world"}) with average degree
  $8$ (\code{DEFAULT\_AVG\_NODE\_DEGREE}) and rewiring probability $0.08$
  (\code{DEFAULT\_REWIRING\_PROB}). Scale-free (Barab\'asi--Albert) and random
  (Erd\H{o}s--R\'enyi) topologies are also implemented as options.
  \item \textbf{Population:} $N = 500$ users plus one platform.
\end{itemize}

\section{Process Overview and Scheduling}
\label{sec:scheduling}

Scheduling is fully deterministic given a seed. At each step, the model's
\code{step()} method executes the following eight phases in strict order;
within phases~1--5 user agents are processed in a fixed list order, so
randomness enters only through the seeded RNG draws inside each submodel. The
model halts when \code{steps} reaches \code{max\_steps}.

\begin{enumerate}
  \item \textbf{Exposure $\rightarrow$ detection $\rightarrow$ harm $\rightarrow$
  trust loss.} For each active user and each active pattern (intensity $>0$):
  a stochastic exposure check (Section~\ref{sub:exposure}); if exposed, a
  detection check (Section~\ref{sub:detection}); then harm
  (Section~\ref{sub:harm}), trust loss (Section~\ref{sub:trustloss}), and harm
  accumulation (Section~\ref{sub:harmgain}). Awareness and positive sentiment
  decay first. After all patterns are processed for a user, per-step caps are
  applied: trust loss is capped at \code{MAX\_TRUST\_LOSS\_PER\_STEP}~$=0.035$
  and harm gain at \code{MAX\_HARM\_GAIN\_PER\_STEP}~$=0.04$.
  \item \textbf{Negative word-of-mouth propagation.} Each active user computes
  its WOM propensity and, if positive, transmits it to up to three network
  neighbours; receivers' staged social signals are then applied to trust
  (Section~\ref{sub:wom}).
  \item \textbf{Recovery via customer support.} Trust partially recovers toward a
  harm-adjusted ceiling, dampened by this-step exposure, cumulative harm, and
  active dark-pattern intensity (Section~\ref{sub:recovery}).
  \item \textbf{Natural passive recovery.} A small additional drift of trust
  toward the harm-adjusted ceiling (Section~\ref{sub:natrecovery}).
  \item \textbf{Positive word-of-mouth.} Only fully satisfied users
  (no exposure this step, trust $\ge 0.60$, zero harm, zero negative WOM, past
  cooldown) spread positive WOM (Section~\ref{sub:poswom}).
  \item \textbf{Churn decisions.} Each active user churns with a logistic
  probability (Section~\ref{sub:churn}); high-harm churners emit a final
  ``exit WOM'' burst. A small background \textbf{natural attrition} then removes
  active users with probability \code{NATURAL\_ATTRITION\_PROBABILITY}~$=0.0001$.
  \item \textbf{Platform update.} The model recomputes churn and the agent-level
  reputation aggregate, then calls \code{PlatformAgent.adapt\_strategy()}
  (Section~\ref{sub:adaptation}), updates economics
  (Section~\ref{sub:economics}), updates the $0$--$100$
  \code{platform\_reputation} (Section~\ref{sub:reputation}), and evaluates the
  four tipping points (Section~\ref{sec:tipping}).
  \item \textbf{Bookkeeping.} Cumulative WOM counters and a ``recent events''
  snapshot are updated, and the \code{DataCollector} records all reporters for
  the step (Section~\ref{sub:observation}).
\end{enumerate}

\section{Design Concepts}
\label{sec:design}

We address each of the eleven ODD design concepts; concepts that do not apply are
stated as such.

\subsection{Basic principles}
The model rests on three principles drawn from the consumer-trust, persuasive-design,
and diffusion literatures: (i) trust is a slowly accumulated, asymmetrically
damaged resource, it is lost faster than it is regained, and accumulated harm
permanently lowers the ceiling to which it can rebound; (ii) the harm a
manipulative interface inflicts depends jointly on the design itself
(detectability, base harm) and on the user (literacy, manipulation sensitivity,
resilience), so detection and harm are heterogeneous; and (iii) dissatisfaction
spreads socially, so individual harm becomes a population-level contagion through
a network. Dark patterns are modelled as an extractive strategy that trades
short-term revenue against long-term trust, support cost, and reputation.

\subsection{Emergence}
The key outcomes, the trajectories of mean trust, negative-WOM prevalence,
cumulative churn, platform reputation, and the revenue/opportunity-cost gap, and
the timing of the four tipping points, \emph{emerge} from the interaction of
individual exposure, detection, trust dynamics, network transmission, and exit
decisions. They are not imposed by any aggregate equation. The differential
fragility of user types likewise emerges from trait differences rather than being
prescribed.

\subsection{Adaptation}
Adaptation is confined to the platform. When \code{adaptive\_platform} is
\code{True}, \code{PlatformAgent.adapt\_strategy()} adjusts the platform's own
behaviour in response to observed churn and reputation: if the per-step churn
rate exceeds \code{CHURN\_ADAPTATION\_THRESHOLD}~$=0.05$, or the agent reputation
drops below $0.55$, or the $0$--$100$ \code{platform\_reputation} drops below
$40$, the platform reduces intensity by
\code{ADAPTATION\_INTENSITY\_REDUCTION}~$=0.08$ and raises support by
\code{ADAPTATION\_SUPPORT\_BOOST}~$=0.03$; if churn is low ($<0.03$) and the
agent reputation is high ($>0.70$), it cautiously raises intensity by
\code{ADAPTATION\_INTENSITY\_INCREASE}~$=0.01$. Users do not adapt their traits;
their only behavioural change is the rise of \code{warning\_awareness} after
receiving negative WOM, which is better described under \emph{Sensing}.

\subsection{Objectives}
The platform's implicit objective, encoded in the adaptation rule, is to sustain
its user base and reputation while extracting revenue: it backs off when churn or
reputation deteriorate and presses harder when they are healthy. User agents have
no explicit objective function; their churn is a probabilistic response to their
state, not an optimisation.

\subsection{Learning}
There is no learning in the reinforcement sense: neither agent updates a policy or
value function from experience, and trait values are fixed. The platform's
adaptation is a fixed reactive rule, not learned. This concept therefore does not
apply.

\subsection{Prediction}
Agents do not form internal predictions of future states. The only forward-looking
quantity is a model-level accounting projection, the counterfactual
no-dark-pattern revenue used to compute \code{opportunity\_cost}, which is a
bookkeeping device, not an agent prediction. This concept does not apply at the
agent level.

\subsection{Sensing}
User agents sense (i) the dark patterns they are exposed to, through the
stochastic detection submodel, whose success rises with \code{digital\_literacy}
and with primed \code{warning\_awareness}; and (ii) signals from network
neighbours, negative WOM (which lowers trust and raises warning awareness) and
positive WOM (which mildly raises trust and positive sentiment). The platform
``senses'' aggregate churn and reputation for its adaptation rule. Sensing is
local: users only perceive their own exposures and their immediate neighbours'
signals.

\subsection{Interaction}
The central interaction is \textbf{word-of-mouth across the social network}.
Negative WOM is a directed, harm-gated, capped transmission from a sender to up
to \code{WOM\_MAX\_NEIGHBORS\_PER\_STEP}~$=3$ active neighbours; its effect on a
receiver is attenuated by a diminishing-returns discount and a trust shield
(Section~\ref{sub:wom}). Positive WOM is a separate, lower-rate transmission from
fully satisfied users. The platform interacts with every user through exposure
and through support-driven recovery. There is no direct user--user interaction
other than WOM.

\subsection{Stochasticity}
Stochasticity enters at four points, all routed through a single seeded RNG
(Mesa's \code{random.Random} seeded via \code{super().\_\_init\_\_(rng=seed)},
with the network builder drawing its own seed from that RNG):
(i) \emph{trait initialization}, user type is drawn from the type
distribution, and each trait is a $\mathrm{Beta}(5,5)$ draw scaled to its range;
(ii) \emph{exposure and detection}, Bernoulli draws against the exposure and
detection probabilities, each perturbed by small uniform noise
(\code{EXPOSURE\_NOISE\_SCALE}, \code{DETECTION\_NOISE\_SCALE}~$=0.05$);
(iii) \emph{WOM transmission}, Bernoulli draws decide whether each
sender--neighbour transmission fires; and (iv) \emph{churn and natural attrition}, Bernoulli draws against the logistic churn probability and the fixed
attrition probability. Because every draw uses the one seeded stream and the
schedule is fixed, a given \texttt{(scenario, seed)} pair is fully reproducible.

\subsection{Collectives}
User types (skeptic, naive, activist) function as \emph{collectives} in the ODD
sense: each is a sub-population whose members share trait-distribution ranges and
whose aggregate trust, churn, and WOM are tracked separately. Membership is fixed
at initialization. The network's small-world community structure is an emergent
topological feature rather than an explicitly represented collective.

\subsection{Observation}
All observation is performed by Mesa's \code{DataCollector}, configured through
\code{build\_all\_reporters()} (Section~\ref{sub:observation}), which records a
fixed set of model-level reporters at the end of every step: aggregate metrics
(mean trust of active and of all users, mean harm, churn rate and cumulative
churn, negative- and positive-WOM counts, warning awareness, the agent and
$0$--$100$ reputations, and the full economic breakdown), per-tipping-point
trigger flags and trigger steps, per-user-type trust/churn/WOM series, and
per-dark-pattern exposure/detection/trust-loss series. The reproduction harness
(\code{reproduce.py}) reads these series to produce all tables and figures.

\section{Initialization}
\label{sec:init}

The model is initialized as follows (all randomness from the seeded RNG):

\begin{enumerate}
  \item \textbf{Dark patterns.} The three \code{DarkPattern} objects are built
  from \code{DARK\_PATTERN\_DEFAULTS}; an enabled pattern receives the global
  \code{dark\_pattern\_intensity}, a disabled one receives intensity $0$.
  \item \textbf{Platform.} A \code{PlatformAgent} is created with the chosen
  intensity, support quality, and adaptivity.
  \item \textbf{Network.} The interaction graph is built (small-world by default;
  the builder draws its own seed from the model RNG, so the topology is
  reproducible).
  \item \textbf{Users.} $N$ \code{UserAgent}s are created and bound to graph
  nodes. Each is assigned a type by a weighted draw from the default
  distribution, \textbf{skeptic $0.30$, naive $0.50$, activist $0.20$}
  (\code{DEFAULT\_TYPE\_DISTRIBUTION}), and each trait is sampled via
  \code{beta\_sample} from that type's range (Table~\ref{tab:traits}) with
  \code{BETA\_SHAPE}~$=5.0$. Each user's \code{trust} and
  \code{perceived\_fairness} are set to its \code{trust\_baseline}; harm, WOM,
  and counters start at zero; \code{active} is \code{True}.
  \item \textbf{Economics and reputation.} The cumulative revenue counter starts
  at \code{INITIAL\_CUMULATIVE\_REVENUE}~$=10{,}000$ (the platform is assumed to
  have prior traction). The initial $0$--$100$ \code{platform\_reputation} is
  drawn \textbf{uniformly from $[50, 70]$} (\code{DEFAULT\_REPUTATION\_RANGE}).

  \emph{Important caveat.} Although each scenario preset in
  \code{config.py} carries a \code{reputation\_range} field, that field is
  \textbf{not consumed by the model constructor}: the helper that expands a
  scenario into constructor keyword arguments does not pass it through, so the
  initial reputation is \emph{always} $U(50,70)$ regardless of the scenario's
  declared range. The \code{reputation\_range} field is therefore presentational
  and has no effect on simulation results.
  \item \textbf{Data collection.} The \code{DataCollector} records the step-0
  state before any dynamics run.
\end{enumerate}

The \emph{seed} is the sole determinant of a run's stochastic realisation; in the
reproduction design seeds $0$--$99$ are used (100 replicates per scenario).

\section{Input Data}
\label{sec:input}

The model uses \textbf{no external input data}. It does not read time series,
empirical traces, GIS layers, or any other external files to drive dynamics. All
quantities are either fixed constants (in \code{config.py}), parameters supplied
at construction (intensity, support, pattern toggles, network and population
settings), or values generated internally by the seeded RNG. Parameter values
were calibrated to literature and behavioural plausibility rather than fitted to a
specific empirical dataset (the calibration rationale is documented separately in
the project's calibration notes).

\section{Submodels}
\label{sec:submodels}

This section gives the exact equations and parameters of every process invoked in
Section~\ref{sec:scheduling}. Unless stated otherwise, all trust-like quantities
are clamped to $[0,1]$ after each update. Symbols: $T$ = trust, $H$ = cumulative
harm, $W$ = negative-WOM propensity, $s_m$ = \code{manipulation\_sensitivity},
$l_d$ = \code{digital\_literacy}, $r_t$ = \code{trust\_resilience},
$SC$ = \code{switching\_cost}.

\subsection{Exposure}
\label{sub:exposure}

For an active user and an active pattern of intensity $I$, exposure is a Bernoulli
event:
\[
\Pr(\text{exposed}) = \max\!\big(0,\; I + \varepsilon\big),
\qquad \varepsilon \sim U(-0.05,\,0.05),
\]
with the noise scale \code{EXPOSURE\_NOISE\_SCALE}~$=0.05$. On exposure the
user's \code{exposure\_count}, \code{cumulative\_exposure} (by $I$), and the
per-pattern counter are incremented.

\subsection{Detection}
\label{sub:detection}

A detection probability is formed from literacy, the pattern's detectability $d$,
and primed awareness $a$ (\code{warning\_awareness}):
\[
p_{\text{detect}} = \min\!\Big(0.95,\;
\max\!\big(0,\; l_d \cdot d \cdot (1 + a) + \varepsilon\big)\Big),
\qquad \varepsilon \sim U(-0.05,\,0.05),
\]
capped at \code{MAX\_DETECTION\_PROBABILITY}~$=0.95$ with noise scale
\code{DETECTION\_NOISE\_SCALE}~$=0.05$. Detection is then a Bernoulli draw against
$p_{\text{detect}}$. Pattern detectabilities are
\code{forced\_trial}~$=0.3$, \code{hard\_cancel}~$=0.4$,
\code{drip\_pricing}~$=0.4$.

\subsection{Harm magnitude}
\label{sub:harm}

The raw harm of a single exposure is
\[
h = b \cdot I \cdot \mu \cdot \sigma_p,
\]
where $b$ is the pattern's \code{base\_harm}
(\code{forced\_trial}~$=1.8$, \code{hard\_cancel}~$=1.6$,
\code{drip\_pricing}~$=2.0$), $\mu$ is \code{detected\_harm\_multiplier} if the
pattern was \emph{detected} (\,$2.2$, $2.0$, $2.0$ respectively\,) or
\code{hidden\_harm\_multiplier} if it was not (\,$0.8$, $0.5$, $0.5$\,), and
$\sigma_p$ is the user's per-pattern \code{pattern\_sensitivity}. The value is
rescaled to $[0,1]$-ish units by dividing by $100$ to give $h_{\text{scaled}}$.

\paragraph{Exposure buildup.}
To model habituation/ramp-in, the first
\code{EXPOSURE\_BUILDUP\_STEPS}~$=3$ encounters with a given pattern deliver only
partial harm. With $k$ the per-pattern exposure count, for $k <3$:
\[
h_{\text{scaled}} \leftarrow h_{\text{scaled}} \cdot
\Big(\,\code{INITIAL\_HARM\_FRACTION} +
(1-\code{INITIAL\_HARM\_FRACTION})\cdot \tfrac{k}{3}\Big),
\qquad \code{INITIAL\_HARM\_FRACTION}=0.2,
\]
so the first hit delivers $20\%$ of harm and full harm applies from the fourth
encounter on.

\subsection{Trust loss from exposure}
\label{sub:trustloss}

The per-exposure trust loss follows the documented formula
\[
\Delta T = -\,\alpha \cdot h_{\text{scaled}} \cdot s_m \cdot (0.6 + 0.8\,l_d)
\cdot (1 - r_t),
\qquad \alpha = \code{ALPHA\_EXPOSURE\_TO\_TRUST} = 0.22 .
\]
The factor $(0.6 + 0.8\,l_d)$ makes more literate users react more strongly to
harm they sustain, and $(1 - r_t)$ lets resilient (mostly naive) users absorb
part of the loss. Perceived fairness drops by $0.8\,|\Delta T|$. After all
patterns are processed in the step, total trust loss is capped at
\code{MAX\_TRUST\_LOSS\_PER\_STEP}~$=0.035$.

\subsection{Harm accumulation}
\label{sub:harmgain}

Cumulative harm grows logistically with a saturation factor, so it slows as
$H\to1$ (modelling desensitisation):
\[
\Delta H = \delta \cdot h_{\text{scaled}} \cdot (0.7 + 0.6\,s_m)
\cdot \max(0,\,1 - H),
\qquad \delta = \code{DELTA\_EXPOSURE\_TO\_HARM} = 0.18 .
\]
Per step, harm gain is capped at \code{MAX\_HARM\_GAIN\_PER\_STEP}~$=0.04$.

\subsection{Negative word-of-mouth}
\label{sub:wom}

\paragraph{Propensity.}
A user generates no negative WOM until accumulated harm passes a cooldown,
\code{WOM\_HARM\_COOLDOWN\_THRESHOLD}~$=0.08$. Past it, a ramp factor
$\rho = \min\!\big(1,\,(H - 0.08)/\code{WOM\_RAMP\_RANGE}\big)$ with
\code{WOM\_RAMP\_RANGE}~$=0.25$ scales a base propensity:
\[
W = \rho \cdot \Big(0.10\,a_s + 0.20\,c_p + 0.45\min(1,H) + 0.25(1 - T)\Big),
\]
where $a_s = $ \code{social\_activity} and $c_p = $ \code{complaint\_propensity}.

\paragraph{Transmission.}
A user with $W>0$ visits its neighbours and stops after at most
\code{WOM\_MAX\_NEIGHBORS\_PER\_STEP}~$=3$ successful contacts. Each contact fires
with probability $a_s \cdot c_p \cdot \code{WOM\_DAMPING\_FACTOR}$
(\code{WOM\_DAMPING\_FACTOR}~$=0.35$). On a firing contact, the receiver's
susceptibility is
\[
\text{impact} = \underbrace{\frac{1}{1 + 0.5\, m}}_{\text{diminishing returns}}
\cdot \underbrace{(1 - 0.60\, T_{\text{nbr}})}_{\text{trust shield}},
\]
with $m$ the number of WOM messages the receiver has already absorbed this step
(\code{WOM\_DIMINISHING\_RATE}~$=0.5$) and \code{WOM\_TRUST\_SHIELD}~$=0.60$. The
receiver's warning awareness rises by
$\code{WOM\_AWARENESS\_BOOST}\cdot(1-a_{\text{nbr}})\cdot\text{impact}$
(\code{WOM\_AWARENESS\_BOOST}~$=0.15$); its trust drops immediately by
$\code{WOM\_TRUST\_PENALTY}\cdot \phi \cdot \text{impact}$ with
\code{WOM\_TRUST\_PENALTY}~$=0.50$ and $\phi=$ \code{social\_influence\_strength}
($0.18$ by default); and a unit social signal $\times\,\text{impact}$ is staged.

\paragraph{Applying the staged signal.}
After all transmissions, each user applies its staged signal to trust:
\[
\Delta T_{\text{social}} = -\,\gamma \cdot S,
\qquad \gamma = \code{GAMMA\_SOCIAL\_TRUST\_LOSS} = 0.12,
\]
where $S$ is the accumulated \code{received\_negative\_signal}; perceived fairness
falls by $0.8\,|\Delta T_{\text{social}}|$, and $S$ resets to $0$. Negative WOM
and positive sentiment also decay each step at rates
\code{NEGATIVE\_WOM\_DECAY\_RATE}~$=0.05$ and
\code{POSITIVE\_WOM\_DECAY\_RATE}~$=0.10$.

\subsection{Recovery via customer support}
\label{sub:recovery}

Trust recovers toward a \emph{harm-adjusted ceiling}
$\bar T = \code{trust\_baseline}\cdot(1-H)$, accumulated harm permanently lowers
how far trust can rebound. The support-driven increment is
\[
\Delta T_{\text{rec}} = \beta \cdot q \cdot d_H \cdot d_E \cdot (1 - I_{\text{plat}}),
\qquad \beta = \code{BETA\_SUPPORT\_RECOVERY} = 0.14,
\]
where $q = $ \code{customer\_support\_quality}; the harm dampening is
$d_H = 1 - \min(H\cdot\code{HARM\_DAMPENING\_FACTOR},\,\code{HARM\_DAMPENING\_CAP})$
with \code{HARM\_DAMPENING\_FACTOR}~$=1.0$ and \code{HARM\_DAMPENING\_CAP}~$=0.85$
(support never falls below $15\%$ effectiveness); the exposure dampening is
$d_E = 1 - \min(1,\, h_{\text{step}}/\code{RECOVERY\_EXPOSURE\_CEILING})$ with
\code{RECOVERY\_EXPOSURE\_CEILING}~$=0.15$ ($h_{\text{step}}$ is total scaled harm
this step); and $I_{\text{plat}}$ is the platform's dark-pattern intensity, so
recovery is suppressed while patterns run. Trust is raised by
$\Delta T_{\text{rec}}$ but not above $\bar T$.

\subsection{Natural passive recovery}
\label{sub:natrecovery}

Independently of support, trust drifts a small amount toward the same ceiling:
\[
\Delta T_{\text{nat}} = \code{NATURAL\_TRUST\_RECOVERY}\cdot(\bar T - T)\cdot(1 - I_{\text{plat}}),
\qquad \code{NATURAL\_TRUST\_RECOVERY}=0.004,
\]
applied only when $T < \bar T$ and again capped at $\bar T$.

\subsection{Positive word-of-mouth}
\label{sub:poswom}

A user spreads positive WOM only if it is fully satisfied, no pattern detected
this step, $T \ge \code{SATISFIED\_TRUST\_THRESHOLD} = 0.60$, zero negative WOM,
zero harm, zero exposure, and past a cooldown of
\code{WOM\_COOLDOWN\_PERIOD}~$=5$ steps since last receiving negative WOM. Each
neighbour contact fires with probability
$\code{POSITIVE\_WOM\_BASE\_RATE}\cdot a_s \cdot T$
(\code{POSITIVE\_WOM\_BASE\_RATE}~$=0.10$); a firing contact raises the
receiver's positive sentiment and nudges its trust up by
$\code{POSITIVE\_WOM\_TRUST\_BOOST}\cdot\phi$
(\code{POSITIVE\_WOM\_TRUST\_BOOST}~$=0.10$), capped at the receiver's
harm-adjusted ceiling.

\subsection{Churn}
\label{sub:churn}

Each active user churns with a logistic probability of an index $z$:
\[
P_{\text{churn}} = \sigma(z),\qquad
z = \theta_0
+ \theta_T \cdot \max\!\big(0,\, (1-T) - D_z\big)
+ \theta_H \cdot H
+ \theta_S \cdot W
- \theta_{SC}\cdot SC
- \text{(retention bonus)} ,
\]
with $\sigma(z) = 1/(1+e^{-z})$ and parameters
$\theta_0 = \code{THETA0} = -8.00$,
$\theta_T = \code{THETA\_TRUST} = 3.50$,
$\theta_H = \code{THETA\_HARM} = 1.90$,
$\theta_S = \code{THETA\_SOCIAL} = 1.20$,
$\theta_{SC} = \code{THETA\_SWITCHING\_COST} = 1.60$. The
\textbf{trust dead-zone} $D_z = \code{CHURN\_TRUST\_DEAD\_ZONE} = 0.30$ is
subtracted from the trust deficit, so trust above $\approx0.70$ contributes
nothing to churn pressure (healthy platforms see only single-digit long-run
churn). Churn is a Bernoulli draw against $P_{\text{churn}}$; on churn the user
becomes inactive, and if its harm is at least
\code{EXIT\_WOM\_HARM\_THRESHOLD}~$=0.30$ it emits one final negative-WOM burst.
Independently, each active user is removed by background \textbf{natural
attrition} with probability \code{NATURAL\_ATTRITION\_PROBABILITY}~$=0.0001$ per
step.

\subsection{Platform economics}
\label{sub:economics}

Let $n_a$ be the number of active users and let the reputation factor be
\[
f_R = \big(\code{platform\_reputation}/100\big)^{\,0.5},
\qquad \text{exponent } \code{REPUTATION\_REVENUE\_EXPONENT} = 0.5 .
\]
Base subscription revenue is
$R_{\text{base}} = n_a \cdot \code{BASE\_REVENUE\_PER\_USER}\cdot f_R$ with
\code{BASE\_REVENUE\_PER\_USER}~$=5.0$. Dark-pattern revenue rewards
\emph{undetected} extraction more heavily than detected extraction: for each
active pattern, with $g$ its \code{short\_term\_gain\_weight}
(\code{forced\_trial}~$=3.0$, \code{hard\_cancel}~$=2.0$,
\code{drip\_pricing}~$=2.5$), $n_{\text{det}}$ detected and $n_{\text{und}}$
undetected exposures this step,
\[
R_{\text{dp}} \mathrel{+}= \Big(n_{\text{det}}\cdot I \cdot g
+ n_{\text{und}}\cdot I \cdot g \cdot \code{HIDDEN\_EXTRACTION\_MULTIPLIER}\Big)\cdot f_R,
\qquad \code{HIDDEN\_EXTRACTION\_MULTIPLIER}=1.5 .
\]
Per-step costs are the sum of churn-replacement cost
($\code{step\_churns}\times\code{CHURN\_REPLACEMENT\_COST}$, $=5.0$), support cost
($n_a\times\code{SUPPORT\_COST\_RATE}\times q$, \code{SUPPORT\_COST\_RATE}~$=0.2$),
and WOM-damage cost
($\code{step\_negative\_wom}\times\code{REPUTATION\_DAMAGE\_COST}$, $=0.5$).
Step profit is revenue minus costs; cumulative revenue, costs, and net value
accumulate over the run. A counterfactual no-dark-pattern projection accumulates a
fixed per-step projected revenue, and \code{opportunity\_cost} is the difference
between cumulative projected and cumulative actual revenue.

\subsection{Platform reputation}
\label{sub:reputation}

The $0$--$100$ \code{platform\_reputation} is updated each step by
\[
R' = \mathrm{clip}\Big(R - p_{\text{churn}} - p_{\text{wom}} + p_{\text{pos}}
+ \code{REPUTATION\_RECOVERY\_RATE},\;
\code{REPUTATION\_FLOOR},\; \code{REPUTATION\_NATURAL\_CAP}\Big),
\]
with penalties and boost expressed per capita and scaled to the $0$--$100$ range:
$p_{\text{churn}} = (\text{step churns}/N)\cdot\code{CHURN\_REPUTATION\_WEIGHT}\cdot100$
($0.8$), $p_{\text{wom}} = (\text{step neg-WOM}/N)\cdot\code{WOM\_REPUTATION\_WEIGHT}\cdot100$
($0.06$), $p_{\text{pos}} = (\text{step pos-WOM}/N)\cdot\code{POSITIVE\_WOM\_REPUTATION\_WEIGHT}\cdot100$
($0.03$); the passive recovery is \code{REPUTATION\_RECOVERY\_RATE}~$=0.10$ per
step; and the value is clipped to $[\code{REPUTATION\_FLOOR},\,
\code{REPUTATION\_NATURAL\_CAP}] = [2.0,\,92.0]$. A separate agent-level
\code{reputation}~$\in[0,1]$ is the convex combination
$0.7\,\overline{T} + 0.3\,(1 - \overline{W})$ of mean active trust and mean active
negative WOM, used in the adaptation rule.

\subsection{Platform adaptation}
\label{sub:adaptation}

See Section~\ref{sec:design} (Adaptation). The rule is a no-op unless
\code{adaptive\_platform} is enabled.

\subsection{Tipping-point detection}
\label{sec:tipping}

After the economic and reputational updates, the model evaluates four conditions
each step. A tipping point is recorded the first time its condition holds for
\code{TIPPING\_POINT\_PERSISTENCE}~$=5$ \emph{consecutive} steps (tracked by a
per-condition streak counter that resets whenever the condition fails). Writing
$\overline{T}$ for mean trust of active users, $\overline{W}$ for their mean
negative WOM, and $C$ for cumulative churn:

\begin{table}[!ht]
\centering
\caption{Tipping-point detection rules (as implemented in
\code{model.\_update\_tipping\_points}). Each must hold for $5$ consecutive steps.}
\label{tab:tipping}
\small
\begin{tabular}{ll}
\toprule
Tipping point & Condition (per step) \\
\midrule
Trust Collapse        & $\overline{T} \le 0.50$ \\
Social Contagion      & $\overline{W} \ge 0.22$ \\
Churn Cascade         & $C \ge 0.35$ \\
Extractive Divergence & revenue gap $\ge 0.20$ of short-term revenue \;\textbf{and}\; $C \ge 0.15$ \\
\bottomrule
\end{tabular}
\end{table}

\noindent Here the revenue gap is
$(\text{short-term} - \text{long-term})/\text{short-term}$, and long-term revenue
is the short-term accumulator discounted by $(1-C)\cdot\overline{T}$. (Note: the
constants \code{TRUST\_COLLAPSE\_THRESHOLD}, \code{CHURN\_ACCELERATION\_THRESHOLD},
and \code{WOM\_BURST\_THRESHOLD} also appear in \code{config.py} but are not the
values consumed by \code{\_update\_tipping\_points}; the operative thresholds are
those in Table~\ref{tab:tipping}, hard-set in \code{model.py}.)

\subsection{Observation submodel}
\label{sub:observation}

At the end of each step the \code{DataCollector} evaluates the reporter dictionary
returned by \code{build\_all\_reporters()}, with the default configuration of
three user types and three patterns this comprises 57 model-level reporters,
spanning aggregate trust/harm/churn/WOM/economic metrics, the four
tipping-point trigger flags and steps, $4\times3$ per-user-type series, and
$3\times3$ per-pattern series. These per-step series are the sole output interface
of the model and are what \code{reproduce.py} aggregates (across 100 seeds, with
$95\%$ normal-approximation confidence intervals) into the manuscript's tables and
figures.

\section*{Reference}
\addcontentsline{toc}{section}{Reference}

\noindent Grimm, Volker, Steven F. Railsback, Christian E. Vincenot, Uta Berger,
Cara Gallagher, Donald L. DeAngelis, Bruce Edmonds, Jiaqi Ge, Jarl Giske,
J\"urgen Groeneveld, Alice S.A. Johnston, Alexander Milles, Jacob Nabe-Nielsen,
J. Gareth Polhill, Viktoriia Radchuk, Marie-Sophie Rohw\"ader, Richard A. Stillman,
Jan C. Thiele, and Daniel Ayllón. 2020. ``The ODD Protocol for Describing
Agent-Based and Other Simulation Models: A Second Update to Improve Clarity,
Replication, and Structural Realism.'' \emph{Journal of Artificial Societies and
Social Simulation} 23 (2): 7. \url{https://doi.org/10.18564/jasss.4259}.

\end{document}
